\section{Future Directions and Open Problems}
\label{sec:future}

The next research frontier is not simply a more capable base model, but a system that can preserve evidence, choose a justified transition, obtain bounded authority, act through changing tools, and verify the effect from authoritative state. The open problems below follow that lifecycle: grounding, long-horizon control, memory, delegation, evaluation, adaptive security, human collaboration, and resource use. Progress should be stated against the object or transition improved (claim, case, proposal, commitment, receipt, or recovery) rather than attributed vaguely to ``agent intelligence.''

\subsection{Grounding as a verifiable contract}

Grounding should connect four contracts: \emph{evidence admission} records whether a source may enter the case; \emph{claim support} links propositions to supporting, contradicting, or indeterminate evidence; \emph{action justification} binds a proposed effect to intent, policy, authority, and current preconditions; and \emph{outcome reconciliation} compares authoritative state with the intended postcondition. Retrieval relevance alone satisfies none of them.

Existing factuality work covers complementary parts of claim support. ALCE separates fluency, correctness, and citation quality~\cite{gao2023alce}; RAGTruth annotates hallucinations in nearly 18,000 RAG responses~\cite{niu2024ragtruth}; FActScore checks atomic facts against a designated source~\cite{min2023factscore}; and VeriScore distinguishes verifiable from unverifiable claims and exposes limited cross-task transfer~\cite{song2024veriscore}. Because their tasks and evaluators differ, the scores should not be pooled. Together they motivate a typed claim ledger that retains evidence edges, contradiction, indeterminacy, and coverage.

Action grounding adds state, authority, and time. ToolSandbox and AppWorld show that dependencies, missing information, and collateral state changes can invalidate an apparently sensible sequence~\cite{lu2025toolsandbox,trivedi2024appworld}; ClawsBench separates useful from unsafe completion, while AgentDojo shows that observations may contain adversarial instructions~\cite{li2026clawsbench,debenedetti2024agentdojo}. A valid proposal must therefore bind admitted observations, exact targets, current preconditions, policy, and decision rights, and completion must be established from an authoritative receipt rather than the model's prediction.

Evaluation should compare RAG, claim ledger, proposal-only, and action-capable systems under the same sources, tools, authority, and budget. The decisive measures are evidence and claim coverage, calibrated abstention, invalid proposals, current preconditions, verified postconditions and collateral effects, lineage through retries or delegation, reviewer corrections, latency, and cost. Stress tests should combine stale or adversarial evidence with changing state and partial execution failure.

The open question is how to make these contracts efficient without letting one model manufacture its own evidence, validation, and authority. Missing support, authority, or postconditions should force retrieval, indeterminacy, review, withheld action, or a reopened case. The objective is not merely fewer unsupported sentences, but fewer unsupported interpretations becoming consequential transitions.

\subsection{Long-horizon control and learned agent policies}

Long-horizon failure is better understood as \emph{state divergence} than as accumulated local mistakes. AgencyBench averages roughly 90 tool calls and one million tokens over 138 tasks~\cite{li2026agencybench}; OSWorld~2.0 reports only 20.6\% binary completion for its strongest system at a 500-step limit on 108 professional workflows~\cite{yuan2026osworld2}. Although the scores are not comparable, both expose lost constraints, stale observations, missed clarification, and skipped verification. More context or more steps does not ensure that open obligations and environment state remain correct.

A long-running agent therefore needs a persistent \emph{execution case} containing acceptance criteria, verified state, open obligations, evidence conflicts, authority, budgets, checkpoints, and receipts. FS-Researcher demonstrates the utility of a durable workspace beyond the context window~\cite{zhu2026fsresearcher}. LongHorizon-Harness goes further by storing task state outside the executor and admitting updates after read-only audit; it improves reported pass rates on WeaveBench, Terminal-Bench~2.1, and OSWorld~2.0~\cite{ma2026longhorizon}. This is comparative, not universal, evidence that independently checked state can matter more than an ever-growing transcript.

Learned policies address distinct parts of the problem: SOAR learns tool consequences, SUPO jointly optimizes tool use and context compression, and AgeMem learns memory operations~\cite{li2026soar,lu2026supo,yu2026agemem}. A safer separation lets the policy propose \emph{what to try next}, while a controller owns evidence admission, invariants, authority, execution, checkpointing, and recovery. The policy may not rewrite the facts, permissions, or success criteria used to judge its transition.

Comparisons among fixed controllers, prompted agents, and learned policies should hold tools, authority, budgets, and stopping rules constant, then measure verified progress, obligation retention, invalid transitions, recovery, postcondition closure, intervention, latency, and cost. Delayed observations, schema changes, interrupted episodes, and validator disagreement should test whether learned compression, planning, memory, or tool choices caused the gain.

The governing invariant is \emph{checkpointed bounded control}: a learned policy proposes from current admitted state, consequential effects remain within an explicit risk and authority envelope, and progress is recorded only from authoritative receipts. The research agenda includes transition-level credit assignment, transfer across tool and policy versions, regression detection, weak-evaluator resistance, and calibrated stopping or escalation. Long-horizon autonomy means preserving justified state transitions, not merely acting for more steps.

\subsection{Memory that is useful, correct, and governable}

The unresolved memory problem is how to govern a stored item's \emph{causal influence} on a later decision. A relevant item may still be false, stale, inferred, cross-user, contradicted, or out of scope. Memory is therefore a lifecycle: admit a typed record, recall it under current scope, judge its applicability, use it in a bounded transition, and confirm, supersede, revoke, or forget it after the outcome.

This lifecycle requires a versioned \emph{memory record} with provenance, time, scope, support, contradiction, and revocation state, plus a \emph{memory use edge} binding that exact version to a claim, parameter, plan step, or proposal. Text2Mem compiles memory instructions into schema-based contracts~\cite{wang2026text2mem}, while Hindsight separates facts, experiences, observations, and opinions~\cite{latimer2026hindsight}. Both improve inspectability, but retrieval performance alone does not establish permissible or causal use in a consequential action.

Benchmarks expose successive gaps. MemBench separates factual from reflective memory and participation from observation~\cite{tan2025membench}; Mem2ActBench shows that seven frameworks remain weak at using memory to choose tools and ground parameters~\cite{shen2026mem2actbench}. Experience-following experiments reveal error propagation through similar retrieved episodes~\cite{xiong2026memorymanagement}, while SteeM shows the tension between useful continuity and anchoring~\cite{huang2026controllable}. The tasks differ, but jointly show that retention, retrieval, application, and reliance control are distinct capabilities.

Learned policies make causal attribution more urgent. AgeMem learns when to store, retrieve, update, summarize, or discard~\cite{yu2026agemem}; Fine-Mem assigns credit to operations that produced downstream evidence and reports gains on Memalpha and MemoryAgentBench~\cite{ma2026finemem}. Retrieval lineage is still not causality: a record may be unused, misleading, or correlated with a favorable evaluator. Interventions that remove, replace, rescope, or revoke individual records should measure the resulting change in claims, proposals, actions, and postconditions, including negative utility.

Memory evaluations should hold the task, model, tools, authority, and budget constant while comparing transcripts, similarity retrieval, structured memory, learned policies, and governed ledgers. They should measure applicability, conflict and revocation handling, correct use in claims and tool arguments, calibrated reliance, collateral effects, cross-user leakage, reviewer burden, latency, and cost under stale, contradictory, poisoned, or deleted records.

The governing invariant is \emph{bounded mnemonic influence}: memory affects a transition only through an inspectable edge from a current, scoped, supported version; it cannot grant authority, certify itself, or survive correction through an untracked derivative. Useful memory is persistent evidence whose influence on bounded action can be inspected, challenged, and revised.

\subsection{Interoperability with trust boundaries}

Dependable interoperability must preserve semantics and responsibility across a trust boundary. MCP standardizes host--tool exchange and A2A agent--agent exchange~\cite{mcp2025spec,a2a2026spec}, but a valid tool description is not permission, a task update is not verified progress, and a successful response is not proof of the intended effect. Standardizing the wire lowers integration cost without removing the local obligation to interpret, authorize, validate, and reconcile what crosses it.

Contract Net and FIPA already formalized task allocation and message semantics~\cite{smith1980contractnet,fipa2002communicative}; modern systems add dynamic discovery, natural language capability claims, black-box services, and consequential tools. MAST identifies 14 failure modes in more than 1,600 multi-agent traces~\cite{cemri2025mast}, while MCP Security Bench finds that stronger instruction following does not imply safer tool composition~\cite{zhang2025mcpsecuritybench}. Although the studies differ, both show why conformant messages and successful components can still produce an invalid global transition.

A delegation should therefore cross the boundary as a typed \emph{delegation envelope}: principal, purpose, inputs and lineage, data scope, capability and effect ceiling, policy, expiry, output contract, and recovery requirements. A linked \emph{result envelope} should bind the producer and version to inputs, transformations, claims, validation, effects, receipts, and open obligations. Together they form a responsibility graph across request, evidence, proposal, authorization, execution, and verification.

Trust admission is a gate, not an endpoint. Attested Tool-Server Admission adds signed clearance, deny-by-default allowlists, and tamper-evident logging to MCP~\cite{metere2026attested}, but cannot establish that a delegated purpose or returned effect is valid for the current case. Discovery should resolve to a locally reviewed, version-pinned manifest; authority and data scope should attenuate at each hop; returned artifacts should re-enter as untrusted evidence; and revocation or schema change should invalidate dependent permissions.

Evaluation should compare raw messages, schema-valid exchange, reviewed manifests, and full delegation contracts under identical agents, tasks, policies, and budgets. It should measure semantic conformance, lineage, authority attenuation, invalid transitions, duplicate effects, verified postconditions, revocation, recovery, accountability, latency, and cost under version drift, replay, substitution, partial artifacts, and false success reports.

The governing invariant is \emph{responsibility-preserving delegation}: every hop retains purpose, disclosed evidence, delegated authority, and the owners of proposal, authorization, commitment, and verification. A downstream hop may narrow purpose, data, duration, or authority, but not broaden them; a returned message may contribute evidence or a proposal, but not silently become permission or proof.

\subsection{Evaluation beyond static leaderboards}

A leaderboard result is a historical join among a benchmark case version, a system manifest, episode records, and an evaluator version. If any element is missing or silently updated, an apparent gain may instead reflect a different scaffold, budget, environment, exposure state, or oracle. Evaluation beyond static leaderboards means maintaining the evidence behind the claim while both system and instrument change.

AgencyBench stresses obligation retention over long, tool-intensive episodes~\cite{li2026agencybench}; Gaia2 uses asynchronous environments and write action verifiers, revealing trade-offs among accuracy, timeliness, and cost~\cite{froger2026gaia2}. Their scores are not comparable, but both make a final answer insufficient: the evaluator must know the observed state, attempted transition, and accepted state change.

The instrument itself also fails. Dynamic construction may reduce some contamination pathways but lacks standardized validity criteria~\cite{chen2025contamination}; search-time leakage can inflate research agent results~\cite{wang2026searchtime}; and the continuing SWE-bench Verified audit shows that curated tasks and executable tests still require inspection of descriptions, coverage, and exposure~\cite{jimenez2024swebench,openai2026swebenchaudit}. Recency alone does not establish validity.

A future benchmark should treat each task as a persistent \emph{benchmark case} binding provenance and time, initial state, constraints, permissions, resource policy, acceptance contract, oracle, adjudication, exposure, health checks, and retirement state. Construction, validation, release, challenge, revision, and retirement should be explicit; corrections must create linked versions rather than silently mixing old and new scores.

No renewal strategy serves every claim: audited cores support longitudinal comparison, held-out waves limit exposure, live environments test changing state, and deployment studies test utility. TRACE evolves tasks while retaining construction trajectories and multi-level validation~\cite{guo2026selfevolving}, but generated difficulty is not automatically realistic difficulty. Portfolio results should remain partitioned by instrument and claim.

The same configured agent should be compared across versioned task waves, tool revisions, policy changes, asynchronous events, and controlled exposure; benchmarks should be challenged with planted defects, collateral effects, oracle omissions, and environment failures. Measures should cover verified transitions, lineage, oracle error, exposure, environment health, repeated success, recovery, rank stability, adjudication, latency, and cost. Generalization claims must remain bounded to sampled variations.

The governing invariant is \emph{versioned evaluative traceability}: every score resolves to immutable cases, a complete system and resource manifest, observed episodes and effects, and evaluator judgments. A leaderboard is then a temporary view over an evaluation ledger, and each result a scoped claim about a named system, instrument, environment, budget, and time.

\subsection{Security for adaptive and connected agents}

Section~\ref{sec:safety} developed security boundaries for evidence, memory, tools, authority, effects, and release. The unresolved problem is whether those boundaries remain valid when the agent changes itself or its composition. A configuration that satisfied an assurance claim at time $t$ may violate it at $t+1$ after a memory consolidation, skill synthesis, policy update, new connector, delegated agent, altered verifier, or learned stopping rule. Adaptation is therefore not simply more data entering the context. It is a transition that can change which evidence is trusted, which behavior is selected, which capabilities are reachable, and which observations count as success.

These changes occur on four coupled planes. The \emph{evidence plane} contains memories, retrieved artifacts, summaries, and learned case state; the \emph{behavior plane} contains prompts, procedures, skills, routing rules, and model or policy parameters; the \emph{capability plane} contains tools, credentials, remote agents, and delegated decision rights; and the \emph{assurance plane} contains policies, monitors, verifiers, review triggers, and rollback criteria. A low-trust trajectory promoted from evidence into a reusable skill crosses from the first plane into the second. If that skill selects a newly installed tool, it reaches the third; if the same learned process also revises the judge that evaluates its outcomes, it weakens the fourth. Security research must preserve these distinctions because provenance on an observation does not authorize a behavior, a useful behavior does not grant a capability, and apparent success does not validate the monitor that declared it.

Recent evidence shows different ways that influence can persist across these planes. BackdoorAgent instruments planning, memory, and tool use attacks across four agent applications and reports, for one GPT-based backbone, trigger persistence of 43.58\%, 77.97\%, and 60.28\% at the respective stages~\cite{feng2026backdooragent}. SkillSec-Eval broadens the unit of analysis to repository admission, semantic retrieval, planner selection, execution, and skill evolution across 327 real-world skills~\cite{badhe2026skillsec}. A very recent preprint on PoisonedEvolution targets the promotion boundary directly: with 10\% attacker support, it reports target behavior embedding in 546 of 600 SkillClaw trials and 369 of 600 trials in a structurally different pipeline~\cite{chen2026poisonedevolution}. Reward misspecification creates a separate endogenous failure: experiments in text-adapted AI Safety Gridworlds find that direct optimization can widen the gap between observed and hidden reward across tested 1.5B--14B language model agents~\cite{cagatan2026rewardhacking}. These studies use different threats, systems, and metrics and should not be pooled. Together they show that cross-stage propagation, poisoned experience promotion, vulnerable skill lifecycles, and proxy optimization are distinct routes by which an apparently useful update can become persistent unsafe control.

Every behavior-affecting change should consequently be represented as a typed \emph{adaptation transaction}. The record should identify the target artifact and exact before--after versions; source trajectories, memories, feedback, reward, and contributors; trust, tenant, purpose, and data scope; the adaptation mechanism and claimed benefit; affected capabilities, policies, invariants, and assurance claims; proposed authority or information flow changes; evaluation and red-team evidence; dependency and derivative artifact links; rollout cohort and effect ceiling; approval, validity, expiry, and revocation state; monitoring thresholds; and rollback or containment procedure. Its lifecycle should distinguish \emph{proposed}, \emph{quarantined}, \emph{evaluated}, \emph{approved}, \emph{canary}, \emph{promoted}, \emph{paused}, \emph{revoked}, and \emph{rolled back}. A model may suggest or synthesize an update, but a separate controller must decide whether that exact version may influence production behavior.

Lineage must survive promotion. If a source trajectory is later found to be adversarial, a user correction is withdrawn, or a tool receipt is invalidated, the system should identify every summary, skill, policy update, evaluator rule, and downstream action that depended on it. Deleting the source text is insufficient when its influence has already been compiled into a derivative artifact. Research is needed on dependency-aware invalidation, counterfactual influence tests, quarantine of suspect derivatives, targeted retraining or replacement, and verification that revocation has reached caches, replicas, active episodes, and delegated agents. When precise removal cannot be established, the safe response is to disable the affected pathway or narrow its authority, not to assume that an opaque update has forgotten its source.

Evaluation should compare a fixed configuration, uncontrolled online adaptation, sandboxed adaptation without promotion controls, and governed adaptation transactions under the same tasks, contributors, tools, authority, and resource budget. Measures should include benign utility gain; hazardous effect escape; update admission precision and recall; cross-stage attack transfer; source-to-derivative lineage coverage; use of revoked or stale derivatives; unauthorized capability expansion; verifier co-adaptation and common-mode failure; canary detection; rollback completeness; residual influence; recovery latency; and cost. Stress tests should mix benign and poisoned trajectories, delay outcome feedback, corrupt reward or evaluators, replay previously successful but now invalid experiences, race policy and artifact versions, revoke a source during an active episode, substitute a remote agent, and distribute a suspect skill across multiple tenants. Hidden safety objectives and independent effect checks are necessary because the adapted component cannot be the sole judge of whether its own adaptation is safe.

The governing invariant should be \emph{non-self-authorizing adaptation}: no memory, skill, learned policy, connector, or evaluator may grant itself trust, broaden executable authority, weaken the controls that govern it, or certify its own successful promotion. Every persistent behavioral change creates a new configured system whose affected assurance claims must be reopened; influence remains bounded by recorded purpose, scope, authority, duration, and effect limits until independent evidence supports a separately authorized promotion. Adaptation can improve how evidence is interpreted and which action is proposed, but it must not silently rewrite the rules that determine which evidence is admissible, which action is permitted, or which receipt closes the case.

\subsection{Human--agent collaboration and accountability}

Section~\ref{sec:safety} specified when a consequential proposal should become a review request. The deeper collaboration problem begins earlier and continues later: how should a shared case allocate evidence, expertise, decision rights, commitments, and recovery responsibilities between people and agents as the work evolves? Maximizing autonomous completion is the wrong objective when the agent lacks relevant context, the human lacks time or technical reach, or only a designated role may accept the consequence. A human being present somewhere in the workflow does not establish meaningful control. The system must represent what each participant knows, may decide, has promised to do, and remains responsible for verifying.

Recent evaluations make this distinction empirical. Collaborative Gym gives humans and agents dual control of a shared environment; in real user comparisons, the strongest collaborative agents were preferred over autonomous counterparts in 86\% of travel planning cases, 74\% of tabular analysis cases, and 66\% of related work cases, yet communication and situational awareness failures still appeared in 65\% and 40\% of real condition trajectories, respectively~\cite{shao2026collaborativegym}. HAS-Bench varies agency level, interaction channel, and participant policy across six domains and finds that human participation can improve both task completion and failure recovery, but that the gain depends on when, how, and by whom the input is exercised~\cite{wu2026hasbench}. A field study of agentic customer service reaches a compatible but domain-specific result: human intervention preserved quality for technical escalations, was less effective for emotional escalations, and supported more post-escalation agent effort when it occurred earlier~\cite{wang2026agentichitl}. These results should not be pooled, but they reject ``human-in-the-loop'' as a sufficient system description. Timing, role, information access, and action authority are causal parts of the configuration.

Collaboration also has a jagged production frontier. In a randomized experiment with 2,234 participants producing 11,024 advertisements, human--AI teams produced 50\% more advertisements per worker and higher-quality text, whereas human--human teams produced higher-quality images; greater delegation to the agent was also associated with more homogeneous outputs~\cite{ju2026collaborating}. Conversely, 1,482 simulated shared workspace sessions show that adding relevant collaborators can lower performance when their contributions are not integrated; shared group memory and designated approval gates improved mean performance most clearly in three-person teams~\cite{kotalwar2026synergy}. One study concerns real creative production and the other simulated data analysis teams, so their effect sizes are not directly comparable. Their common lesson is structural: complementary capability does not become joint performance unless responsibility and evidence are routed to the decision where they matter.

An ongoing collaboration should therefore be governed by a typed \emph{collaboration contract}, not reconstructed from chat history. This record should bind the shared goal and acceptance criteria to participant identities and roles; evidence each participant supplied, can access, or is expected to obtain; delegated tasks and human-reserved decisions; open questions, disagreements, and uncertainty; current artifact and proposal versions; communication, clarification, handoff, takeover, and escalation triggers; deadlines and safe defaults; effect ceilings; and the owners of authorization, execution, verification, appeal, and recovery. Its lifecycle should distinguish \emph{negotiated}, \emph{accepted}, \emph{active}, \emph{waiting}, \emph{handed off}, \emph{proposed}, \emph{authorized}, \emph{committed}, \emph{verified}, \emph{disputed}, and \emph{repaired}. Unlike a review request, which asks whether one exact proposal may cross a boundary, the contract preserves the evolving division of labor that produces, authorizes, and checks that proposal.

Human input must be typed as carefully as agent output. A clarification, fact attestation, preference, artifact edit, policy interpretation, one-case exception, action approval, and reusable instruction have different scopes and should create different state transitions. A person's correction of one paragraph must not silently become durable memory; approval of one action must not validate the evidence behind it or authorize a later variant; and an agent's incorporation of feedback must remain linked to the exact input and resulting artifact. The case ledger should make it possible to reconstruct who could see which evidence, who held the relevant decision right, which proposal was authorized, what effect was committed, who checked the receipt, and who can contest or repair the outcome. Explanations are then decision-specific views over this record, not substitutes for it.

Evaluation should compare autonomous execution, fixed approval gates, learned deferral, an unstructured shared workspace, and typed collaboration contracts under the same tasks, evidence, authority, models, tools, and budgets. Outcome measures should be joined with missed and unnecessary handoffs; handoff timing; clarification value; evidence and obligation transfer completeness; decision right violations; rework; takeover and recovery success; human workload and situational awareness; calibration; output diversity; appeal accessibility; latency; and cost. Stress tests should withhold expertise from one participant, introduce conflicting evidence and role changes, delay human responses, interrupt a handoff, revise an artifact after approval, make an approver unavailable, distribute knowledge across several people, and require recovery after an agent or human error. Studies should also report the participant population and role competence rather than treating a simulated user, crowd worker, domain expert, and accountable decision owner as interchangeable.

The governing invariant should be \emph{accountable joint control}: no participant may make a commitment beyond recorded role and authority, and every transfer of work must preserve the material evidence, open obligations, decision rights, and recovery owner needed for the next participant to act. Human approval authorizes only the exact proposal reviewed; agent completion neither erases human responsibility nor manufactures it. A human is meaningfully in the loop only when the relevant evidence reaches them early enough, they possess the decision right at issue, and they can change, stop, contest, or repair the resulting transition. The research target is not maximum autonomy or maximum intervention, but a collaboration whose changing division of labor remains explicit, bounded, and auditable from evidence to effect.

\subsection{Economics and sustainability}

Section~\ref{sec:evaluation} treated resources as part of the episode record. The unresolved control problem is whether an agent can decide, before resources are spent, which additional retrieval, model call, tool, verifier, collaborator, or retry is justified by the evidence still missing and the consequence at stake. Cost is therefore not merely a bill attached to a completed trajectory. It is a constraint on which transition may be attempted next. Monetary spend, tokens, compute, wall-clock time, tool quotas, human attention, energy, and emissions belong in the same typed control state as quality, authority, and risk, while remaining distinct quantities that cannot be collapsed into one price.

Current evidence shows why raw capability comparisons are economically incomplete. AgencyBench reports averages of roughly 90 tool calls and one million tokens per task across long, feedback-rich episodes~\cite{li2026agencybench}. In a different setting, compute-matched evaluation of reasoning strategies finds that simpler chain-of-thought self-consistency frequently outperforms more elaborate methods and that multi-agent debate or Reflexion can become worse as additional compute is supplied~\cite{wang2024tokeneconomies}. These results use different tasks, models, and resource definitions and should not be compared numerically. Together they show that more steps or tokens do not constitute progress unless the extra expenditure improves a verified outcome, reduces relevant uncertainty, or lowers downstream risk.

Budget awareness is itself an agent capability, not a deployment wrapper. EcoAgent-Bench assigns priced actions and explicit budgets to 304 real-derived tasks that require agents to decide whether to use local evidence, broaden search, select a stronger model, escalate to a human, or stop on an unsupported premise. Across seven evaluated agents, tool API configurations achieved only 3.9--24.0\% micro strict success and at most 7.3\% economic consistency, with common failures divided between premature stopping and overspending on inexpensive cases~\cite{wu2026ecoagentbench}. Conversely, LLM-as-Scheduler shows that a structured gate and dynamic workflow routing can reduce tokens by 43\% and end-to-end latency by more than 36\%, with at most a 1.4 percentage-point accuracy reduction relative to its strong fixed workflow baseline~\cite{xiang2026scheduler}. The first study tests whether an agent chooses economically; the second tests one architecture for avoiding unnecessary workflow stages. Their combination motivates controllers that can spend selectively without interpreting a large budget as permission or a small budget as justification for unverified completion.

Sustainability requires an equally explicit measurement boundary. An initial empirical and theoretical study of web agents finds that architectural choices materially change energy and carbon estimates, that higher energy consumption does not necessarily produce better results, and that undisclosed model and process details limit credible accounting~\cite{krupp2025sustainablewebagents}. A token count alone cannot establish energy use because hardware, batching, model implementation, datacenter region, electricity mix, tool services, and time all matter. Nor does a per-episode efficiency gain establish an aggregate reduction if lower cost induces substantially more episodes. Claims about energy, carbon, or water should therefore identify the measured or estimated components, hardware and provider boundary, workload and concurrency, region and time, allocation method, uncertainty, and whether induced demand is inside the claim.

Each episode should consequently carry a versioned \emph{resource contract}. The contract should state the goal and minimum acceptable evidence, quality, safety, and effect verification; monetary ceiling and currency; token, compute, tool, storage, latency, and human attention limits; pricing and carbon intensity sources with versions and timestamps; per-stage reservations and contingency; reusable artifact and cache policies; permitted model, tool, and human escalation tiers; marginal value and stopping rules; and the principal authorized to amend the budget. A linked \emph{resource ledger} should record, for every transition, predicted, reserved, and realized consumption; evidence gained or uncertainty reduced; artifact and effect identifiers; retries, waste, reuse, and cache provenance; current balance; forecast cost-to-completion; and any price, latency, or emissions estimate revision. Its lifecycle should distinguish \emph{proposed}, \emph{admitted}, \emph{reserved}, \emph{spent}, \emph{reconciled}, \emph{reforecast}, \emph{exhausted}, and \emph{closed}. This makes resource exhaustion a typed state that can trigger a narrower proposal, a justified amendment, an incomplete artifact, or escalation rather than an improvised low-quality answer.

The accounting boundary must extend beyond inference. Total cost of ownership includes evidence acquisition, indexing and storage, tool and connector fees, human clarification and review, evaluation and red teaming, incident response, compensation and rollback, monitoring, benchmark refresh, policy and connector maintenance, and retirement or migration of persistent records. Savings from a cheaper model can be erased by more retries, review corrections, or escaped defects; an additional verifier can increase direct cost while reducing expected incident loss. Economic claims should therefore compare the agentic system with an appropriate deterministic workflow, fixed RAG or automation pipeline, assistive human workflow, and fully manual baseline on the same demand, acceptance criteria, authority, and accounting horizon. The relevant output is net value and a quality--cost--latency--risk frontier, not cost per model call in isolation.

Evaluation should compare fixed budgets, adaptive budgets with explicit amendment, static workflow routing, dynamic routing, and unconstrained execution under the same task distribution and system manifest. Measures should include verified effects per resource envelope; cost and energy to accepted outcome; marginal evidence or quality gain per transition; budget forecast calibration; premature stopping and unnecessary escalation; overspend and stranded reservations; reuse validity; human time; tail latency; incident and recovery cost; maintenance burden; and the fraction of environmental impact that is measured rather than imputed. Stress tests should change model and tool prices mid-episode, exhaust a quota, delay the inexpensive model, invalidate a cached plan, raise carbon intensity, introduce an apparently cheap tool with hidden review cost, and present a rare high-risk case near budget exhaustion. Hard quality, authority, safety, and verification requirements must remain gates rather than soft terms that an optimizer can trade away.

The governing invariant should be \emph{resource-bounded sufficiency}: every expenditure must be attributable to an authorized objective and a versioned resource contract, and every claim of savings must remain joined to the outcome quality, risk, and accounting boundary under which it was obtained. The system may route, reuse, stop, or escalate to avoid waste, but it may not economize by admitting weaker evidence, bypassing required validation, broadening authority, hiding human work, or declaring an unverified effect complete. Sustainable agency is not minimum computation at any cost; it is enough justified computation, human attention, and external service use to transform admissible evidence into a bounded and verified effect, followed by a defensible stop.

\section{Conclusion}
\label{sec:conclusion}

This survey has used \emph{LLM-centered agentic systems} as the umbrella for a systems family, not a synonym for autonomous chatbots. Autoregression describes token generation; it does not specify evidence, persistent state, authority, recovery, or proof of completion. Agency emerges at the system boundary where a model adaptively helps choose the next action while tools, policies, state, evaluators, and people constrain the choice. The central transition is from generating a response to transforming evidence into a bounded, reviewable, and verified effect over time.

RAG, workflows, bounded agents, and multi-agent systems are alternative execution forms, not a maturity ladder. RAG interfaces with evidence; workflows specify predictable transformations; bounded agents select among permitted transitions when the useful next step cannot be fixed; and multi-agent systems distribute work across additional boundaries. ReAct and Toolformer show how models can interleave generation and tool interaction, but tool choice alone supplies neither durable state nor permission nor effect verification~\cite{yao2023react,schick2023toolformer}. A more adaptive form is justified only when it resolves uncertainty a simpler form cannot.

The forms are connected by a persistent case whose typed records survive model calls, tools, participants, and time. Evidence records support claims; scoped memories preserve experience; plans expose dependencies; proposals remain distinct from authorizations; action contracts bind permitted effects; receipts distinguish acknowledgement from verified change; and review, delegation, governance, and resource records preserve responsibility. This chain prevents recurring category errors: retrieved text is not automatically admissible evidence, memory is not permission, a plausible plan is not an authorized action, and a successful API response is not proof of the intended effect.

Comparative evaluation supports this view. GAIA, WebArena, OSWorld, SWE-bench, and AgencyBench test different combinations of information gathering, interface control, software change, state, and tools~\cite{mialon2024gaia,zhou2024webarena,xie2024osworld,jimenez2024swebench,li2026agencybench}; AppWorld and $\tau$-bench add state change and repeated run checks~\cite{trivedi2024appworld,yao2024taubench}. Because tasks, scaffolds, budgets, environments, and evaluators differ, scores cannot be pooled into a universal measure of agency. The common lesson is narrower: fluent behavior and component success do not establish a valid evidence-to-effect trajectory.

Deployment evidence similarly favors bounded augmentation. Assistance improved productivity in customer support and software development studies while people retained review and release responsibility~\cite{brynjolfsson2025generativeaiatwork,cui2026highskilledwork}. Human--AI advertising teams produced more output and higher-quality text, whereas human--human teams produced higher-quality images; greater AI delegation also increased homogeneity~\cite{ju2026collaborating}. These outcomes are not one effect size, but they show why deployment claims must retain the baseline, authority, human work, accepted outcome, and downstream correction.

Dependable composition therefore follows proposal--gate--execute--verify. A model may interpret evidence and propose a claim, plan, tool call, memory update, or adaptation; a governed controller checks schema, provenance, scope, policy, budget, and preconditions before trusted state changes. Execution returns a receipt reconciled with the intended postcondition, and uncertain validation triggers repair, rollback, escalation, or an explicitly incomplete state. Delegation and adaptation may extend capability, but must not silently expand authority.

The credible path forward is disciplined, evidence-bearing composition: use the simplest execution form that can meet the goal; make state, authority, uncertainty, and obligations explicit; treat model outputs as proposals rather than facts or permissions; verify effects against the world; retain enough lineage to repair failure; and stop when evidence, authority, or resources no longer support another action. A dependable agentic system is not the one that acts most often, but the one that can show why an action was justified, who authorized it, what changed, whether the intended effect holds, and when not to act.
